% chapters/16/latihan.tex
\bagian{Latihan dan Refleksi}

\begin{enumerate}
  \item Jelaskan mengapa pembuat permainan tidak dapat menilai karyanya
        sendiri dengan baik. Berikan satu contoh dari pengalaman Anda sendiri
        selama mengerjakan proyek akhir.

  \item Sebuah level memiliki $v_0 = 700$ piksel per detik, $g = 1800$ piksel
        per detik kuadrat, dan kecepatan jalan 300 piksel per detik. Petaknya
        berukuran 32 piksel. Hitunglah tinggi lompatan tertinggi dan jarak
        mendatar terjauh pada ketinggian yang sama.

  \item Mengapa jarak mendatar terjauh bergantung pada beda tingginya?
        Apakah melompat ke tempat yang lebih rendah memberi jarak yang lebih
        jauh atau lebih dekat? Jelaskan alasannya.

  \item Pada peta ketiga contoh keluaran, pemeriksaan keterjangkauan
        mengembalikan 15,6 persen sedangkan simulasinya mengembalikan 82,0
        persen. Jelaskan mengapa kedua angka itu berbeda, dan mengapa angka
        yang lebih tinggi justru lebih menyesatkan.

  \item Sebutkan sekurang-kurangnya empat hal yang tidak dapat diperiksa oleh
        pemeriksaan otomatis pada bab ini.

  \item Ubah satu pernyataan GDD berikut menjadi butir \istilah{checklist}
        yang dapat diperiksa: ``Pemain harus merasa terburu-buru.''

  \item Sebuah catatan \istilah{playtest} berbunyi ``permainannya agak
        membingungkan''. Ubahlah catatan itu menjadi catatan yang dapat
        ditindaklanjuti, dan jelaskan keterangan apa yang perlu ditambahkan.

  \item Tiga temuan berikut harus ditimbang beratnya. Urutkan dari yang paling
        berat, dan jelaskan alasan setiap penempatan.
        \begin{enumerate}
          \item Warna latar pada satu bagian level salah.
          \item Pemain dapat terjatuh ke dalam lubang dan tidak dapat keluar.
          \item Bunyi kalah tidak terdengar ketika nyawa habis.
        \end{enumerate}

  \item Jelaskan perbedaan antara ``GDD diperbarui'' dan ``belum dikerjakan''
        pada Tabel~\ref{tab:gdd-build}. Mengapa perbedaan itu penting bagi
        penilaian?
\end{enumerate}

\begin{refleksibox}
\begin{enumerate}
  \item Pada sesi \istilah{playtest}, temuan mana yang paling mengejutkan
        bagi kelompok Anda? Apakah temuan itu muncul pada pemeriksaan
        otomatis juga?

  \item Apakah Anda sempat ingin memberi petunjuk kepada penguji? Bagian mana
        yang membuat Anda ingin menolongnya, dan apa yang akan Anda perbaiki
        pada level Anda sebagai akibatnya?

  \item Setelah menyelesaikan seluruh materi mata kuliah ini, bagian mana yang
        paling Anda kuasai, dan bagian mana yang masih perlu Anda pelajari?
        Tuliskan satu langkah nyata yang akan Anda lakukan untuk yang kedua.
\end{enumerate}
\end{refleksibox}
